\section{Noyau topologique de phase et génération du catalogue de germes}
\label{act21:tpk:capitulo}

Les deux évaluations de l'APP acquièrent une efficacité générative lorsque l'on précise
quelle évaluation agit, en quelle position elle s'effectue, comment cette position se déplace
et quelle information chaque transition conserve. Le \emph{Topological Phase Kernel},
abrégé en TPK et dénommé en espagnol \emph{núcleo topológico de fase}, réunit
ces opérations dans une dynamique sur des états enrichis. Sa définition
comprend la sélection, le transport, la mise à jour arithmétique et la conservation du
registre. La construction finie développée dans ce chapitre fournit
les germes sur lesquels agiront les relèvements régionaux et le prolongement
coinductif.

Le résultat principal de cette étape est un catalogue généré et classé :
les conditions initiales des deux curseurs produisent des émissions entières ;
leur lecture ternaire détermine des mots orientés ; l'action diédrale organise
ces mots en orbites. Chaque étape conserve ses préimages et les relations
qui permettent de revenir à la description précédente. La chaîne de cardinaux
\[
 104\,976\longrightarrow468\longrightarrow243\subset729,
 \qquad 243\longrightarrow43,
\]
désigne sous une forme abrégée des objets différents. Chaque objet est ensuite construit et chaque dénombrement
est justifié, avant que l'un d'eux soit utilisé comme donnée de
la sélection régionale.

\subsection{Sélection de phase, transport et état enrichi}
\label{act21:tpk:operaciones}

\subsubsection{Sélecteur tritique de phase opérationnelle}

La phase nonadique prend ses valeurs en \(D_9=\{1,\ldots,9\}\). Le sélecteur  
tritique est la fonction
\begin{equation}
 M_{\mathrm{ph}}:D_9\longrightarrow\{-1,0,+1\},\qquad
 M_{\mathrm{ph}}(g)=
 \begin{cases}
 +1,&g\in\{1,4,7\},\\
 -1,&g\in\{2,5,8\},\\
 0,&g\in\{3,6,9\}.
 \end{cases}
 \label{act21:tpk:selector}
\end{equation}
Dans la réalisation émettrice, la valeur \(+1\) active l'évaluation additive,
\(-1\) active l'évaluation multiplicative et \(0\) enregistre l'événement
neutre. L'orientation cardinale du curseur est une autre coordonnée de l'état.
Cette séparation permet de fixer simultanément quelle opération est lue et dans quelle
direction est transporté son point d'évaluation.

Pour les transitions numérotées par \(t\geq1\), la phase précédant la lecture est
\begin{equation}
 g_t=1+[(t-1)]_9,\qquad \epsilon_t=M_{\mathrm{ph}}(g_t),
 \label{act21:tpk:fase}
\end{equation}
où \([n]_b\) désigne le représentant de \(n\) modulo \(b\) dans
\(\{0,\ldots,b-1\}\). Une fenêtre de neuf transitions contient la suite
\((+1,-1,0,+1,-1,0,+1,-1,0)\).

\subsubsection{Curseurs orientés et transport cardinal}

Des indices \(I_9=\{0,\ldots,8\}\) sont utilisés pour les positions et des étiquettes positives \((a,b)=(i+1,j+1)\) pour les évaluations APP. Un curseur fini est
\[
 c=(i,j;d)\in\mathcal S:=I_9^2\times\mathcal D,
 \qquad \mathcal D=\{N,E,S,O\}.
\]
Les quatre directions agissent par
\[
 v_N=(-1,0),\quad v_E=(0,1),\quad
 v_S=(1,0),\quad v_O=(0,-1).
\]
L'application cardinale correspondant à \(d\) est
\[
 T_d(i,j)=\bigl([i+(v_d)_1]_9,[j+(v_d)_2]_9\bigr).
\]
Ainsi, \(M_{\mathrm{ph}}\) et \(T_d\) ont des domaines, des codomaines et des actions différents. L'activation tritique sélectionne l'un des deux curseurs ; le transport cardinal modifie sa position dans la direction que ce curseur conserve déjà.

Le relèvement entier du curseur est
\(\widetilde c=(x,y;d)\in\mathbb Z^2\times\mathcal D\), avec position finie
\(([x]_9,[y]_9)\). Ses nombres de tours sont récupérés en
\begin{equation}
 x=9\lfloor x/9\rfloor+[x]_9,\qquad
 y=9\lfloor y/9\rfloor+[y]_9.
 \label{act21:tpk:vueltas}
\end{equation}
Ces deux divisions concernent le transport spatial. La division d'une
évaluation additive ou multiplicative constitue un registre arithmétique
distinct, qui sera conservé conjointement avec elles.

\subsubsection{Composition de la transition}

L'état de la réalisation finie comprend les deux curseurs relevés,
leurs directions, la phase, les accumulateurs de la fenêtre active, la liste
ordonnée des blocs déjà émis et le registre des transitions. Chaque registre
local contient les positions antérieures, les évaluations entières des deux
feuillets, leurs restes et quotients, le régime actif et les orientations.
Les champs de préfixe et de frontière qui apparaîtront dans un prolongement seront
incorporés au même état avec leurs règles de compatibilité.

La transition s'ordonne comme
\begin{equation}
 \mathcal U_t=\mathsf{Mem}_t\circ\mathsf{Tra}_t\circ\mathsf{Sel}_t.
 \label{act21:tpk:composicion}
\end{equation}
La sélection détermine le régime et conserve l'échantillon antérieur à tout déplacement. Le transport applique le mouvement actif et l'inversion cardinale prescrite par le calendrier. L'actualisation incorpore cet échantillon aux accumulateurs et au registre ; à la fin d'une fenêtre, elle émet son bloc. La composition utilise donc un échantillon antérieur au transport bien que son archivage définitif s'effectue à la fin de la transition.

\subsection{Conditions initiales et énumération réversible}
\label{act21:tpk:semillas}

Les feuillets additif et multiplicatif disposent de curseurs indépendants. Le
domaine des conditions initiales est le produit ordonné
\[
 \Omega_{\mathrm{seed}}=\mathcal S_+\times\mathcal S_\times,
 \qquad |\mathcal S|=9^2\cdot4=324.
\]
Par conséquent,
\begin{equation}
 |\Omega_{\mathrm{seed}}|=324^2=104\,976.
 \label{act21:tpk:censo-inicial}
\end{equation}
Le nombre \(81\) compte des positions ; \(324\) compte des positions orientées d'un curseur ; \(104\,976\) compte des paires ordonnées de curseurs. Le choix d'une position en APP conserve uniquement une partie d'une condition initiale de l'émetteur.

On fixe la numérotation
\[
 \operatorname{ind}(N)=0,\quad\operatorname{ind}(E)=1,\quad
 \operatorname{ind}(S)=2,\quad\operatorname{ind}(O)=3,
\]
et l'on définit
\begin{equation}
 \nu(i,j;d)=4(9i+j)+\operatorname{ind}(d).
 \label{act21:tpk:indice-semilla}
\end{equation}

\begin{proposition}[Énumération exacte des conditions initiales]
\label{act21:tpk:enumeracion}
L'application \(\nu\) est une bijection de \(\mathcal S\) sur
\(\{0,\ldots,323\}\). Le couple d'indices se code de manière réversible
par \(324\nu_++\nu_\times\in\{0,\ldots,104\,975\}\).
\end{proposition}
\begin{proof}
La division par quatre restaure
\[
 \operatorname{ind}(d)=[\nu]_4,\qquad
 j=[\lfloor\nu/4\rfloor]_9,\qquad i=\lfloor\nu/36\rfloor.
\]
Les bornes du domaine garantissent que ces trois quantités appartiennent à
leurs ensembles d'indices et reconstruisent \(\nu\) au moyen de
\eqref{act21:tpk:indice-semilla}. Pour le couple, la division euclidienne par
\(324\) restitue le quotient \(\nu_+\) et le reste \(\nu_\times\).
\end{proof}

L'énumération parcourt le domaine complet avant tout regroupement
régional. Les indices servent à localiser les préimages d'une
émission, en conservant explicitement la condition qui a produit chaque trajectoire.

\subsection{Évaluation des feuillets et registre arithmétique}
\label{act21:tpk:lecturas}

En une position positive \((a,b)\), les évaluations entières sont
\(A_+(a,b)=a+b\) et \(A_\times(a,b)=ab\). Pour \(n>0\), on conserve
\begin{equation}
 \rho_9(n)=1+[n-1]_9,\qquad
 q_9(n)=\lfloor(n-1)/9\rfloor,\qquad
 n=\rho_9(n)+9q_9(n).
 \label{act21:tpk:division-evaluacion}
\end{equation}
L'activation additive incorpore \(\rho_9(A_+)\) à l'accumulateur correspondant. L'activation multiplicative utilise un observable du résidu \(\rho_9(A_\times)\), défini ci-dessous.

\subsubsection{Exposant discret dans les unités modulo neuf}

Les puissances de \(2\) modulo \(9\) parcourent
\[
 2^0,2^1,\ldots,2^5\equiv1,2,4,8,7,5\pmod9,
 \qquad 2^6\equiv1\pmod9.
\]
Le logarithme discret dans ce groupe de six unités se représente par
\begin{equation}
 \ell_9(1)=0,\quad\ell_9(2)=1,\quad\ell_9(4)=2,
 \quad\ell_9(8)=3,\quad\ell_9(7)=4,\quad\ell_9(5)=5.
 \label{act21:tpk:logaritmo}
\end{equation}
La loi émettrice étend l'observable par le biais de \(\ell_9(3)=\ell_9(6)=\ell_9(9)=0\). Pour conserver son origine, l'échantillon enregistré inclut le résidu original et l'indicateur d'appartenance au groupe d'unités. Ainsi, une valeur observée nulle peut provenir de l'unité \(1\) ou de l'un des trois résidus non inversibles, et son origine demeure récupérable.

% BEGIN R27_MG05_CALIBRES
% MG05a. Tras el párrafo de extensión y procedencia que sigue a
% act21:tpk:logaritmo; antes de «Orden de lectura y movimiento».
\subsubsection{Jauges de l'observable multiplicatif et normalisation du générateur}
\label{mg05:calibres}

Le choix de l'exposant discret admet une comparaison finie dans laquelle restent fixes le support APP, les conditions initiales des deux curseurs, le sélecteur tritique, le calendrier d'inversion et la règle émettrice. La variation se concentre sur l'observable du feuillet multiplicatif. Soient \(\ell_{2}\) et \(\ell_{5}\) les logarithmes du groupe \((\mathbb Z/9\mathbb Z)^\times\) de générateurs \(2\) et \(5\), respectivement, prolongés par zéro sur \(\{3,6,9\}\). Ils sont comparés à \(v_3\), évaluée sur les représentants positifs \(1,\ldots,9\), et à la lecture identité sur ces mêmes représentants.

\begin{proposition}[Normalisation du logarithme multiplicatif]
\label{mg05:normalizacion}
Les deux isomorphismes du groupe des unités modulo neuf avec
\(\mathbb Z/6\mathbb Z\), normalisés en envoyant un générateur sur
\(1\), sont \(\ell_2\) et \(\ell_5\). Ils satisfont
\[
 \ell_5(x)\equiv-\ell_2(x)\pmod6
 \qquad\bigl(x\in(\mathbb Z/9\mathbb Z)^\times\bigr).
\]
Le choix du plus petit générateur positif fixe \(\ell_2=\ell_9\),
avec les valeurs de \eqref{act21:tpk:logaritmo}.
\end{proposition}
\begin{proof}
Les puissances de \(2\) parcourent \(1,2,4,8,7,5\) et se referment à la
sixième étape. Les générateurs d'un groupe cyclique d'ordre six
correspondent aux exposants premiers avec six : \(1\) et \(5\).
Leurs représentants sont \(2\) et \(2^5\equiv5\pmod9\).
Comme \(5\equiv2^{-1}\pmod9\), exprimer une unité comme puissance de
\(5\) change le signe de son exposant en base \(2\). La
normalisation par le plus petit représentant positif sélectionne \(2\).
\end{proof}

La condition d'homomorphisme seule autorise également des applications
d'image propre. La proposition distingue donc la représentation
fidèle du groupe des unités de ses quotients. Le prolongement par zéro
détermine en outre le traitement des trois résidus non inversibles ;
le registre de provenance conserve celui d'entre eux qui a été observé.

La comparaison des quatre observables parcourt les mêmes \(104\,976\) conditions initiales et regroupe leurs émissions de six fenêtres. En désignant par \(E_g\) l'émetteur d'observable multiplicatif \(g\), on obtient le recensement d'images suivant.
\begin{table}[htbp]
\centering
\small
\begin{tabular}{@{}lrrrl@{}}
\toprule
Observable & \(|\operatorname{im}E_g|\) & Racines nulles
 & Ley aditiva & Normalisation\\
\midrule
\(\ell_2\)&468&oui&oui&plus petit générateur\\
\(\ell_5\)&468&oui&oui&generador inverso\\
\(v_3\)&144&no&oui&image unitaire nulle\\
\(\operatorname{id}\)&684&non&non&résidu positif\\
\bottomrule
\end{tabular}
\caption{Comparaison de quatre jauges sur un même transport et un même domaine de conditions initiales.}
\label{mg05:tabla-calibres}
\end{table}

Les deux logarithmes satisfont la loi additive des exposants. La valuation
\(v_3\) vaut zéro sur toutes les unités, mais prend les valeurs
\(1,1,2\) sur \(3,6,9\) ; d'où la différence entre leurs deux colonnes
de contrôle. L'identité conserve le résidu positif et ne satisfait pas la loi
d'homomorphisme vers \(\mathbb Z/6\mathbb Z\) ; par exemple, pour
l'unité \(1\), l'égalité exigée serait \(1\equiv2\pmod6\).
Le recensement de \(468\) émissions est commun aux deux logarithmes et
conserve donc l'ambiguïté d'inversion. Le critère du plus petit
générateur résout cette ambiguïté avant toute lecture régionale
de constantes. Les cardinaux du tableau décrivent l'évaluation
exhaustive de ces quatre observables sur le domaine fixé ; leur
fonction est de contrôler la normalisation opérationnelle de l'émetteur.

% END R27_MG05_CALIBRES

\subsubsection{Ordre de lecture et mouvement}

Les évaluations s'effectuent aux positions antérieures à la transition.
Si \(\epsilon_t=+1\), on accumule la lecture additive et on déplace son curseur.
Si \(\epsilon_t=-1\), on accumule l'exposant discret multiplicatif et on
déplace le second curseur. Si \(\epsilon_t=0\), on incrémente le compteur
neutre et on conserve les deux positions.

À la fin des étapes \(27\) et \(54\), les deux directions sont remplacées
par leurs opposées. Les fenêtres s'enchaînent en conservant leurs positions et
orientations. Au début de chaque fenêtre, seuls ses trois accumulateurs locaux sont remis à zéro ; le registre précédent reste disponible.

\begin{proposition}[Inversion du transport relevé]
\label{act21:tpk:inversion}
Soient \(a,b\in\{0,1\}\) les indicateurs d'activation et d'inversion d'un
curseur. Si \(\operatorname{opp}\) échange \(N\leftrightarrow S\) et
\(E\leftrightarrow O\), la transformation
\[
 F_{a,b}(z,d)=\bigl(z+a v_d,\operatorname{opp}^{\,b}(d)\bigr)
\]
admet pour inverse
\[
 F_{a,b}^{-1}(z',d')=
 \bigl(z'-a v_{\operatorname{opp}^{\,b}(d')},
       \operatorname{opp}^{\,b}(d')\bigr).
\]
La réduction spatiale modulo neuf entrelace cette transformation avec le
transport fini.
\end{proposition}
\begin{proof}
Appliquer deux fois l'opposition restaure la direction initiale. Une fois
cette direction récupérée, la soustraction de \(a v_d\) annule le déplacement.
Pour la projection spatiale, il suffit de l'égalité
\([z+a v_d]_9=[[z]_9+a v_d]_9\), composante par composante. L'inversion de
direction conserve cette égalité.
\end{proof}

La phase détermine les indicateurs utilisés par l'inverse. Le registre
spatial et l'arithmétique peuvent ainsi être conservés conjointement, même en
traversant le bord du représentant \(I_9^2\).

\subsection{Six fenêtres et émission décimale par blocs}
\label{act21:tpk:emision}

\subsubsection{Accumulation dans une fenêtre nonadique}

La fenêtre \(m\), avec \(1\leq m\leq6\), comprend les étapes
\(9m-8,\ldots,9m\). Soient \(S_m\) la somme de ses résidus positifs
additifs actifs, \(E_m\) la somme de ses exposants discrets actifs et
\(Z_m\) son nombre d'événements neutres.

\begin{lemma}[Multiplicités d'activation]
\label{act21:tpk:tres-eventos}
Chaque fenêtre contient trois activations additives, trois multiplicatives et
trois neutres. En particulier, \(Z_m=3\).
\end{lemma}
\begin{proof}
Les phases d'une fenêtre parcourent exactement \(1,\ldots,9\). Chacun
des trois ensembles de \eqref{act21:tpk:selector} a pour cardinal trois.
\end{proof}

Les six fenêtres forment \(54\) transitions et produisent six blocs.  
Cette longueur appartient à l'émetteur initial défini ici. La continuation  
en profondeur agit ensuite sur les mots et leurs registres compatibles.

\subsubsection{Règle de codification décimale}

La loi émettrice spécifie les chiffres
\begin{align}
 d_{m,1}&=[S_m]_{10},\notag\\
 d_{m,2}&=[S_m+E_m]_{10},\notag\\
 d_{m,3}&=[3S_m+5E_m+7Z_m]_{10},\notag\\
 U_m&=100d_{m,1}+10d_{m,2}+d_{m,3}.
 \label{act21:tpk:bloque-decimal}
\end{align}
Les poids \(100,10,1\) sont les positions des centaines, des dizaines et des unités.
L'indice \(10\) indique le reste décimal, et chaque bloc satisfait
\(0\leq U_m\leq999\). Le mot émis est la suite ordonnée
\(\mathbf U=(U_1,\ldots,U_6)\), de largeur fixe de trois chiffres par bloc.

Les coefficients \(3,5,7\) constituent les poids explicites de cette loi
émettrice. Une fois la règle fixée, chaque chiffre se calcule à partir des
accumulateurs, et les résultats de ce chapitre se déduisent de cette règle.
Le nombre \(1\) qui apparaît sous sa forme réduite a une origine
additionnelle concrète : \(7Z_m=21\equiv1\pmod{10}\).

Si \(s_m=[S_m]_{10}\) et \(e_m=[E_m]_{10}\), on obtient
\begin{equation}
 G(s_m,e_m)=100s_m+10[s_m+e_m]_{10}
                   +[3s_m+5e_m+1]_{10}=U_m.
 \label{act21:tpk:acoplamiento-firmas}
\end{equation}
La lettre \(e_m\) désigne ici une composante de la signature d'exposants discrets. Sa définition appartient à l'émetteur fini.

\begin{proposition}[Récupération des deux signatures]
\label{act21:tpk:recuperacion-firmas}
L'application \(G:\{0,\ldots,9\}^2\to\{0,\ldots,999\}\) est injective.
Pour un bloc \(U=100a+10b+c\) de son image, l'inverse est
\[
 s=a,\qquad e=[b-a]_{10},\qquad
 c=[3s+5e+1]_{10}.
\]
En conséquence, le mot \(\mathbf U\) détermine les deux signatures de six composantes \(\mathbf s\) et \(\mathbf e\).
\end{proposition}
\begin{proof}
Les deux derniers termes de \eqref{act21:tpk:acoplamiento-firmas} appartiennent à
\(\{0,\ldots,99\}\), de sorte que la centaine récupère \(s\). La dizaine
est \([s+e]_{10}\) ; en soustrayant \(s\) et en réduisant modulo dix, on récupère
\(e\). L'unité vérifie la troisième congruence. La reconstruction s'applique
indépendamment aux six positions du mot.
\end{proof}

Les totaux entiers \(S_m,E_m\) se retrouvent à partir des lectures enregistrées ;
les centaines et les dizaines retrouvent leurs restes. Par exemple, un total
\(S_m=15\) produit \(s_m=5\). Cette différence identifie exactement
quelle information chaque niveau de codage conserve.

% BEGIN R27_DEC01
\subsubsection{Alphabet décimal et congruence de contrôle de l'émetteur}
\label{sec:rec27-codigo}

On conserve la loi émettrice de \eqref{act21:tpk:bloque-decimal},
avec les lectures APP et l'observable \(\ell_9\) de
\eqref{act21:tpk:logaritmo}, les curseurs de
\ref{act21:tpk:operaciones} et le calendrier de six fenêtres de
\ref{act21:tpk:emision}. Les accumulateurs sont \(S_m,E_m,Z_m\), avec
\(Z_m=3\), et les chiffres de sortie sont notés \(d_{m,1},d_{m,2},d_{m,3}\).
La proposition~\ref{act21:tpk:recuperacion-firmas} restitue déjà les deux
signatures résiduelles. Leur combinaison avec les images effectives des
accumulateurs détermine l'alphabet suivant.

\begin{proposition}[Congruence de contrôle du troisième chiffre]
Chaque bloc émis satisfait
\[
 d_{m,3}=[5d_{m,2}-2d_{m,1}+1]_{10}.
\]
\end{proposition}
\begin{proof}
L'inverse de la proposition~\ref{act21:tpk:recuperacion-firmas}
donne \([E_m]_{10}=[d_{m,2}-d_{m,1}]_{10}\). En substituant
$E_m\equiv d_{m,2}-d_{m,1}\pmod{10}$ et $Z_m=3$ dans le troisième chiffre, on obtient
$3d_{m,1}+5(d_{m,2}-d_{m,1})+21\equiv5d_{m,2}-2d_{m,1}+1\pmod{10}$.
\end{proof}

\begin{proposition}[Alphabet décimal de quarante-deux blocs]
L'image de l'émetteur par fenêtre est constituée exactement des blocs
\[
\begin{array}{c|rrrrrrr}
[E_m]_{10}&\multicolumn{7}{c}{U_m}\\ \hline
0&114&227&443&556&669&885&998\\
1&129&232&458&561&674&890&903\\
3&149&252&478&581&694&810&923\\
5&169&272&498&501&614&830&943\\
7&189&292&418&521&634&850&963\\
9&109&212&438&541&654&870&983
\end{array}
\]
\end{proposition}
\begin{proof}
Une trajectoire cardinale lit trois termes consécutifs du cycle additif
\[
 (2,3,4,5,6,7,8,9,1).
\]
Leurs sommes sont
$(9,12,15,18,21,24,18,12,6)$ ; inverser la direction conserve l'ensemble
des sommes. Par conséquent
\[
 \operatorname{Im}S_m=\{6,9,12,15,18,21,24\},\qquad
 \operatorname{Im}[S_m]_{10}=\{1,2,4,5,6,8,9\}.
\]
Dans le canal multiplicatif, un facteur fixe divisible par trois apporte zéro. Pour les facteurs inversibles, les sommes cycliques de trois lectures sont
\[
\begin{array}{c|c}
\text{facteur fixe}&\text{sommes possibles}\\\hline
1,8&\{1,3,7,9\}\\
2,7&\{3,5,9\}\\
4,5&\{1,5,7\}
\end{array}
\]
Il en résulte $\operatorname{Im}E_m=\{0,1,3,5,7,9\}$.
L'atteignabilité peut déjà être vérifiée dans la première fenêtre : tous les
curseurs des deux tableaux suivants ont pour direction $E$.
\[
\begin{array}{c|rrrrrrr}
S_1&6&9&12&15&18&21&24\\\hline
(i_+,j_+)&(0,8)&(0,0)&(0,1)&(0,2)&(0,3)&(0,4)&(0,5)
\end{array}
\]
\[
\begin{array}{c|rrrrrr}
E_1&0&1&3&5&7&9\\\hline
(i_\times,j_\times)&(2,0)&(0,0)&(0,1)&(1,1)&(0,2)&(0,4)
\end{array}
\]
L'indépendance des deux curseurs réalise chacun des $7\cdot6$
couples. Les deux premiers chiffres restituent le couple résiduel, de sorte que ses
42 images sont distinctes. La congruence de contrôle fixe le troisième chiffre
et produit exactement le tableau énoncé.
\end{proof}

Le bloc décimal conserve ainsi les deux signatures résiduelles au moyen d'un codage injectif et d'un chiffre de contrôle. Ce résultat décrit l'alphabet de l'émetteur fini ; le prolongement des lectures régionales conserve en outre les conditions de phase, d'orientation et de mémoire de l'état enrichi.

L'image par fenêtre comporte donc \(42\) blocs. La concaténation
compatible de six fenêtres conserve le recensement de \(468\) émissions de
la proposition~\ref{act21:tpk:468}, dont la réduction ternaire produit
\(243\) mots dans l'espace ambiant de \(729\). Chaque cardinal compte
l'objet de son étape et conserve la provenance du même émetteur.

% END R27_DEC01

\subsubsection{Codification positionnelle du mot de six blocs}

Les six blocs admettent une codification entière à largeur fixe,
\begin{equation}
 \mathscr N(\mathbf U)=\sum_{m=1}^{6}U_m1000^{6-m},
 \qquad0\leq\mathscr N(\mathbf U)<1000^6.
 \label{act21:tpk:entero-base1000}
\end{equation}
Le premier bloc occupe la position de plus haut poids et le sixième, celle des unités de base mille. Un bloc \(058\) occupe une position complète avec une valeur entière \(58\) ; sa largeur de trois chiffres fait partie de la présentation décimale de cette position.

\begin{proposition}[Récupération positionnelle des émissions]
\label{act21:tpk:inversa-base1000}
La codification \eqref{act21:tpk:entero-base1000} détermine le mot complet
par le biais de
\[
 U_m=\left[\left\lfloor
 \frac{\mathscr N(\mathbf U)}{1000^{6-m}}\right\rfloor\right]_{1000},
 \qquad1\leq m\leq6.
\]
\end{proposition}
\begin{proof}
En divisant par \(1000^{6-m}\), les blocs précédents produisent des multiples entiers de \(1000\), le bloc \(m\) produit \(U_m\), et la contribution des suivants appartient à \([0,1)\). La partie entière élimine cette dernière contribution et le résidu modulo \(1000\) élimine les premières.
\end{proof}

Cette codification conserve les six émissions. Son domaine possède des positions ordonnées de base mille, tandis que le domaine APP possède des positions spatiales modulo neuf. La composition maintient les deux indices : \(m\) identifie une fenêtre d'émission et \((i,j)\) identifie la cellule visitée lors d'une transition de cette fenêtre.

\subsection{Calcul complet d'une émission qui commence par 501}
\label{act21:tpk:ejemplo501}

On fixe les conditions initiales
\[
 c_{+,0}=(0,4;N),\qquad c_{\times,0}=(2,3;N),
 \qquad \nu_+=16,\quad\nu_\times=84.
\]
Les positions positives initiales sont \((1,5)\) et \((3,4)\). L'indice
chronologique précédent est \(t=0\), la première phase lue est \(g_1=1\) et le
registre est vide. Ces coordonnées sont les entrées
du calcul ; les blocs s'obtiennent en exécutant la récurrence.

\begin{table}[htbp]
\caption{Première fenêtre : lectures précédentes et accumulation active.}
\label{act21:tpk:tabla-501}
\small
\begin{tabular}{@{}rcccrrr@{}}
\toprule
\(t\)&\(\epsilon_t\)&position \(+\)&position \(\times\)&
\(\Delta S\)&\(\Delta E\)&\((S,E,Z)\)\\
\midrule
1&\(+1\)&\((1,5)\)&\((3,4)\)&6&0&\((6,0,0)\)\\
2&\(-1\)&\((9,5)\)&\((3,4)\)&0&0&\((6,0,0)\)\\
3&0&\((9,5)\)&\((2,4)\)&0&0&\((6,0,1)\)\\
4&\(+1\)&\((9,5)\)&\((2,4)\)&5&0&\((11,0,1)\)\\
5&\(-1\)&\((8,5)\)&\((2,4)\)&0&3&\((11,3,1)\)\\
6&0&\((8,5)\)&\((1,4)\)&0&0&\((11,3,2)\)\\
7&\(+1\)&\((8,5)\)&\((1,4)\)&4&0&\((15,3,2)\)\\
8&\(-1\)&\((7,5)\)&\((1,4)\)&0&2&\((15,5,2)\)\\
9&0&\((7,5)\)&\((9,4)\)&0&0&\((15,5,3)\)\\
\bottomrule
\end{tabular}
\notatabla{Les positions sont des étiquettes positives des cellules antérieures au mouvement. Les événements neutres augmentent uniquement \(Z\).}
\end{table}

Les trois évaluations additives actives sont
\[
 1+5=6,\qquad9+5=14=5+9,\qquad8+5=13=4+9,
\]
de telle sorte que \(S_1=6+5+4=15\). Les trois évaluations multiplicatives actives
sont \(3\cdot4=12\), \(2\cdot4=8\) et \(1\cdot4=4\). Leurs résidus
positifs sont \(3,8,4\), avec des observables \(0,3,2\) ; par conséquent,
\(E_1=5\). Enfin \(Z_1=3\). La règle décimale donne
\[
 d_{1,1}=[15]_{10}=5,\quad
 d_{1,2}=[15+5]_{10}=0,\quad
 d_{1,3}=[45+25+21]_{10}=1,
\]
et on obtient \(U_1=501\).

La seconde lecture additive précédente a pour quotient arithmétique \(1\),
tandis que son curseur relevé se trouve en \((-1,4;N)\), avec un nombre de tours
verticaux \(-1\). Ce sont deux données différentes du même registre. Toutes deux
interviennent dans la récupération intégrale de l'évaluation et de sa position.

\begin{table}[htbp]
\caption{Les six fenêtres de la condition initiale \((16,84)\).}
\label{act21:tpk:tabla-seis501}
\begin{tabular}{@{}rrrrrrrr@{}}
\toprule
\(m\)&pas&\(S_m\)&\(E_m\)&\(Z_m\)&\(s_m\)&\(U_m\)&\([-U_m]_3\)\\
\midrule
1&1--9&15&5&3&5&501&0\\
2&10--18&6&5&3&6&614&1\\
3&19--27&24&5&3&4&498&0\\
4&28--36&21&5&3&1&169&2\\
5&37--45&12&5&3&2&272&1\\
6&46--54&12&5&3&2&272&1\\
\bottomrule
\end{tabular}
\end{table}

Le mot résultant et ses signatures sont
\begin{align*}
 \mathbf U&=(501,614,498,169,272,272),\\
 \mathbf s&=(5,6,4,1,2,2),&
 \mathbf e&=(5,5,5,5,5,5).
\end{align*}
L'inversion cardinale après l'étape \(27\) modifie les parcours des
trois fenêtres finales. Le bloc répété \(272\) correspond à deux
fenêtres successives possédant leurs propres positions et leur propre registre chronologique.

Comme deuxième cas, la condition initiale minimale \((\nu_+,\nu_\times)=(0,0)\)
produit
\[
 \mathbf U=(252,109,252,903,850,850),\quad
 \mathbf s=(2,1,2,9,8,8),\quad
 \mathbf e=(3,9,3,1,7,7).
\]
Dans ce cas, les deux curseurs retournent à leur position et orientation initiales
après \(54\) transitions. La phase nonadique effectue six tours, le
compteur chronologique avance de \(54\) unités et le registre contient
\(54\) échantillons. La lecture positionnelle du retour et
l'histoire enregistrée décrivent des propriétés distinctes de l'exécution.

\subsection{Factorisation du recensement fini}
\label{act21:tpk:censo}

Les signatures \(f_+(\nu)=\mathbf s\) et \(f_\times(\nu)=\mathbf e\) se
calculent en exécutant le curseur correspondant pendant les six fenêtres.
L'indépendance des positions et des directions entre les deux curseurs donne
\begin{equation}
 T_{\mathrm{em}}(\nu_+,\nu_\times)
   =G^{(6)}\bigl(f_+(\nu_+),f_\times(\nu_\times)\bigr),
 \label{act21:tpk:factorizacion}
\end{equation}
où \(G^{(6)}\) applique \eqref{act21:tpk:acoplamiento-firmas} à chaque position.
La factorisation permet d'énumérer d'abord \(324+324\) trajectoires et
de combiner ensuite les signatures distinctes, en conservant la multiplicité de
leurs préimages.

\par\noindent\begin{minipage}{\linewidth}
\subsubsection{Images des deux signatures}

Les tableaux suivants comprennent toutes les signatures. Leur multiplicité est le
nombre de conditions initiales qui les produisent ; l'indice minimal localise
l'une de ces conditions. La préimage complète est définie par
\[
 f_\sigma^{-1}(a)=\{\nu\in\{0,\ldots,323\}:f_\sigma(\nu)=a\}.
\]
La récurrence des sections précédentes détermine chaque élément de cet
ensemble par un calcul entier fini.

\begin{table}[H]
\centering\fontsize{9.5}{11.4}\selectfont
\setlength{\abovecaptionskip}{8pt}\setlength{\belowcaptionskip}{6pt}
\renewcommand{\arraystretch}{1.12}
\caption{Les dix-huit signatures additives.}
\label{act21:tpk:firmas-aditivas}
\begin{tabular}{@{}lrr@{\qquad}lrr@{}}
\toprule
signature&multiplicité&\(\nu_{\min}\)&signature&multiplicité&\(\nu_{\min}\)\\
\midrule
\texttt{122564}&18&17&\texttt{564122}&18&16\\
\texttt{122898}&18&24&\texttt{645212}&18&4\\
\texttt{212645}&18&5&\texttt{654889}&18&33\\
\texttt{212988}&18&0&\texttt{889221}&18&13\\
\texttt{221456}&18&29&\texttt{889654}&18&32\\
\texttt{221889}&18&12&\texttt{898122}&18&25\\
\texttt{456221}&18&28&\texttt{898465}&18&20\\
\texttt{465898}&18&21&\texttt{988212}&18&1\\
\texttt{546988}&18&9&\texttt{988546}&18&8\\
\bottomrule
\end{tabular}

\par\vspace{12pt}

\caption{Les vingt-six signatures multiplicatives.}
\label{act21:tpk:firmas-multiplicativas}
\begin{tabular}{@{}lrr@{\qquad}lrr@{}}
\toprule
signature&multiplicité&\(\nu_{\min}\)&signature&multiplicité&\(\nu_{\min}\)\\
\midrule
\texttt{000000}&108&8&\texttt{555555}&24&84\\
\texttt{177393}&8&1&\texttt{555717}&8&14\\
\texttt{177555}&8&54&\texttt{555771}&8&18\\
\texttt{177771}&8&11&\texttt{555933}&8&24\\
\texttt{339555}&8&6&\texttt{717555}&8&12\\
\texttt{339771}&8&15&\texttt{717717}&8&21\\
\texttt{339933}&8&59&\texttt{717933}&8&19\\
\texttt{393177}&8&0&\texttt{771177}&8&9\\
\texttt{393393}&8&45&\texttt{771339}&8&13\\
\texttt{393555}&8&40&\texttt{771555}&8&16\\
\texttt{555177}&8&52&\texttt{933339}&8&57\\
\texttt{555339}&8&4&\texttt{933555}&8&26\\
\texttt{555393}&8&41&\texttt{933717}&8&17\\
\bottomrule
\end{tabular}
\end{table}
\end{minipage}\par


\begin{proposition}[Image et multiplicités de l'émetteur]
\label{act21:tpk:468}
L'image \(\mathcal U_6=\operatorname{im}T_{\mathrm{em}}\) contient
\(468\) mots distincts. Ses préimages se répartissent en \(432\)
fibres de cardinal \(144\), \(18\) de cardinal \(432\) et \(18\) de
cardinal \(1944\).
\end{proposition}
\begin{proof}
L'évaluation exhaustive de l'algorithme de la section ~\ref{act21:tpk:algoritmo} sur les indices \(0,\ldots,323\) donne les deux tableaux ci-dessus. La première partition a \(18\) classes de cardinal \(18\), et la seconde \(24\) classes de cardinal \(8\), une de cardinal \(24\) et une de cardinal \(108\). Les sommes \(18\cdot18=324\) et \(24\cdot8+24+108=324\) vérifient la masse totale des deux partitions. La proposition ~\ref{act21:tpk:recuperacion-firmas} rend injectif le couplage de ses images ; par conséquent, il y a \(18\cdot26=468\) émissions. Par \eqref{act21:tpk:factorizacion}, la préimage d'un couple de signatures est le produit de ses deux préimages. On obtient \(18\cdot24=432\) produits de cardinal \(18\cdot8=144\), \(18\) de cardinal \(18\cdot24=432\) et \(18\) de cardinal \(18\cdot108=1944\). Enfin,
\[
 432\cdot144+18\cdot432+18\cdot1944=104\,976.
\]
La partie énumérative est une vérification finie exacte de la récurrence ;
la factorisation explique pourquoi ce décompte couvre toutes les paires.
\end{proof}

L'émission de l'exemple \ref{act21:tpk:ejemplo501} a une fibre de cardinal
\(18\cdot24=432\). L'exemple d'indices \((0,0)\) a une fibre de
cardinal \(18\cdot8=144\). Dans les deux cas, le mot émis récupère
les signatures, tandis que l'identification d'une condition initiale particulière
utilise son indice ou son registre de trajectoire.

\subsection{Lecture ternaire orientée et conservation des fibres}
\label{act21:tpk:reduccion-ternaria}

La réduction orientée agit séparément sur chaque bloc :
\begin{equation}
 \Psi:\mathcal U_6\longrightarrow\mathbb F_3^6,\qquad
 \Psi(U_1,\ldots,U_6)=([-U_1]_3,\ldots,[-U_6]_3).
 \label{act21:tpk:psi}
\end{equation}
Son signe fixe une orientation du catalogue. L'identification équilibrée
\(0\mapsto0\), \(1\mapsto+1\), \(2\mapsto-1\) s'applique ensuite,
composante par composante.

Le module neuf organise les évaluations et positions d'APP ; les restes modulo dix forment chaque chiffre ; la base mille regroupe trois positions décimales ; le module trois détermine la signature ternaire de chaque bloc. Ces quatre usages ont des objets et des fonctions définis. En particulier, \(\Psi\) lit six blocs et produit six trits ; une conversion positionnelle de l'entier formé par dix-huit chiffres serait une autre opération.

Pour les deux conditions calculées, on obtient
\begin{align*}
 \Psi(501,614,498,169,272,272)&=(0,1,0,2,1,1),\\
 \Psi(252,109,252,903,850,850)&=(0,2,0,0,2,2).
\end{align*}
L'émetteur décimal est défini avant cette réduction. La fonction de
\(\Psi\) est de produire l'objet ternaire sur lequel agissent la symétrie
orbitale et les sélecteurs régionaux. Cette fonction explique sa position dans la
généalogie : elle relie une émission déjà calculée à une classification structurale.

L'image \(\mathcal W_6:=\Psi(\mathcal U_6)\) contient exactement
\(243\) mots au sein de l'ensemble ambiant de \(3^6=729\) mots.
L'énumération de l'algorithme final donne l'histogramme
\begin{equation}
\begin{array}{c|rrrrrr}
 |\Psi^{-1}(w)|&1&2&3&4&5&6\\\hline
 \text{nombre de mots }w&132&66&18&3&6&18.
\end{array}
\label{act21:tpk:fibras-ternarias}
\end{equation}
Les sommes des deux lectures sont
\[
132+66+18+3+6+18=243,
\qquad132+2\cdot66+3\cdot18+4\cdot3+5\cdot6+6\cdot18=468.
\]
La première compte les mots ternaires ; la seconde récupère le nombre d'émissions réparties entre ses fibres.

\begin{example}[Deux émissions avec la même signature ternaire]
\label{act21:tpk:colision}
Les émissions
\[
 (498,501,614,252,252,109),\qquad
 (870,810,923,252,252,109)
\]
sont distinctes et toutes deux ont pour image \((0,0,1,0,0,2)\) par \(\Psi\).
Chacune a \(144\) conditions initiales. Le mot ternaire conserve
une classe de lecture ; les émissions et leurs préimages distinguent les données
qui convergent vers cette classe.
\end{example}

La conservation des fibres s'exprime concrètement par
\[
 \mathcal F_w=\{\mathbf U\in\mathcal U_6:\Psi(\mathbf U)=w\},
 \qquad
 \Omega_{\mathbf U}=T_{\mathrm{em}}^{-1}(\mathbf U).
\]
Ainsi, la préimage complète de \(w\) dans le domaine initial est l'union disjointe des ensembles \(\Omega_{\mathbf U}\), avec \(\mathbf U\in\mathcal F_w\). Le registre enrichi peut conserver l'émission, ses signatures et la condition initiale ainsi que sa lecture ternaire. L'inverse exacte se réfère alors à la donnée enregistrée avec son type ; le mot de six trits pris isolément conserve la portée précise de \eqref{act21:tpk:psi}.

\subsection{Action diédrale et les quarante-trois orbites}
\label{act21:tpk:orbitas}

Sur six coordonnées se définissent les permutations
\begin{align}
 r(u_1,u_2,u_3,u_4,u_5,u_6)
   &=(u_3,u_1,u_2,u_5,u_6,u_4),\\
 s(u_1,u_2,u_3,u_4,u_5,u_6)
   &=(u_4,u_5,u_6,u_1,u_2,u_3).
 \label{act21:tpk:generadores-diedrales}
\end{align}
La première fait tourner les deux groupes de trois coordonnées dans des sens conjugués ;
la seconde échange les deux groupes. Leur composition satisfait
\(r^3=s^2=1\) et \(srs=r^{-1}\). Elles engendrent ainsi une action du groupe diédral
\(D_3\), à six éléments.

\begin{proposition}[Équivariance et stabilité du catalogue]
\label{act21:tpk:equivarianza}
La réduction \(\Psi\) commute avec l'action de \(D_3\). Les ensembles
\(\mathcal U_6\) et \(\mathcal W_6\) sont stables sous leurs générateurs.
\end{proposition}
\begin{proof}
L'égalité \(\Psi(g\mathbf U)=g\Psi(\mathbf U)\) se vérifie en chaque
coordonnée parce que \(g\) permute les positions et \(\Psi\) applique la même
réduction à chaque position. La stabilité du catalogue constitue une autre
vérification : en générant ses \(468\) éléments par
\eqref{act21:tpk:factorizacion}, on vérifie que \(r\mathbf U\) et
\(s\mathbf U\) appartiennent à nouveau à cet ensemble. L'algorithme de la
section~\ref{act21:tpk:algoritmo} effectue ces vérifications d'appartenance.
L'équivariance transmet ensuite la stabilité à son image ternaire.
\end{proof}

\begin{proposition}[Décomposition orbitale complète]
\label{act21:tpk:43}
L'ensemble \(\mathcal W_6\) se décompose en \(43\) orbites : \(38\)
de cardinal six et \(5\) de cardinal trois.
\end{proposition}
\begin{proof}
Pour chaque \(w\in\mathcal W_6\) on forme
\[
 \mathcal O(w)=\{w,rw,r^2w,sw,srw,sr^2w\}.
\]
On choisit comme représentant le mot lexicographiquement minimal de cet ensemble. Deux mots ont le même représentant exactement lorsqu'ils appartiennent à la même orbite. L'algorithme génère la liste des représentants du tableau~\ref{act21:tpk:representantes} ; la somme de leurs cardinaux vaut \(38\cdot6+5\cdot3=243\). Comme chaque mot du catalogue participe à cette opération et que le catalogue est stable, la liste couvre toute l'image.
\end{proof}

\begin{table}[htbp]
\caption{Représentants minimaux et cardinaux des quarante-trois orbites.}
\label{act21:tpk:representantes}
\small
\begin{tabular}{@{}lr@{\qquad}lr@{\qquad}lr@{}}
\toprule
représentant&cardinal&représentant&cardinal&représentant&cardinal\\
\midrule
\texttt{001002}&6&\texttt{002112}&6&\texttt{012222}&6\\
\texttt{001011}&6&\texttt{002211}&6&\texttt{021022}&6\\
\texttt{001012}&6&\texttt{002220}&6&\texttt{021121}&6\\
\texttt{001021}&6&\texttt{011021}&6&\texttt{021202}&6\\
\texttt{001022}&6&\texttt{011112}&6&\texttt{022022}&3\\
\texttt{001100}&3&\texttt{011120}&6&\texttt{022112}&6\\
\texttt{001101}&6&\texttt{011201}&6&\texttt{022121}&6\\
\texttt{001110}&6&\texttt{011211}&6&\texttt{022202}&3\\
\texttt{001112}&6&\texttt{011220}&6&\texttt{022211}&6\\
\texttt{001202}&6&\texttt{011222}&6&\texttt{022222}&6\\
\texttt{001210}&6&\texttt{012112}&6&\texttt{112112}&3\\
\texttt{001211}&6&\texttt{012121}&6&\texttt{112211}&3\\
\texttt{001220}&6&\texttt{012202}&6&\texttt{112222}&6\\
\texttt{001222}&6&\texttt{012210}&6&&\\
\texttt{002012}&6&\texttt{012220}&6&&\\
\bottomrule
\end{tabular}
\end{table}

Les orbites héritent des invariants entiers de leurs fibres. Pour chaque mot,
\(n(w)=|\mathcal F_w|\) compte les émissions, tandis que
\[
 M(w)=\sum_{\mathbf U\in\mathcal F_w}|\Omega_{\mathbf U}|
\]
compte les conditions initiales. Le recensement des symboles
\(c(w)=(n_0,n_1,n_2)\) enregistre combien de fois chaque classe ternaire apparaît.
Ces données, conjointement avec l'orientation et le support des parcours,
fournissent les prédicats de la sélection régionale ultérieure. La
classification conserve ainsi plus d'information que la forme d'un mot
représentant.

\subsection{Algorithme autonome d'énumération et de contrôle}
\label{act21:tpk:algoritmo}

L'énumération utilise les entiers, les listes finies et les opérations déjà définies.  
Le procédé suivant précise son exécution ; \(\sigma\) prend les valeurs  
\(+\) et \(\times\).

\begin{enumerate}
\item Pour chaque \(\nu=0,\ldots,323\), récupérer \((i,j;d)\) par
l'inverse de \eqref{act21:tpk:indice-semilla}. Initialiser une liste vide de signatures.
\item Pour chaque fenêtre \(m=1,\ldots,6\), initialiser l'accumulateur \(A=0\).
Parcourir \(t=9m-8,\ldots,9m\), calculant \(\epsilon_t\) par
\eqref{act21:tpk:fase}.
\item Si \(\sigma=+\) et \(\epsilon_t=+1\), ajouter
\(\rho_9(i+j+2)\) à \(A\) et mettre à jour
\((i,j)\leftarrow T_d(i,j)\). Si \(\sigma=\times\) et
\(\epsilon_t=-1\), ajouter
\(\ell_9(\rho_9((i+1)(j+1)))\) et effectuer le même transport.
\item À la fin de \(t=27\) ou \(t=54\), remplacer \(d\) par son opposée.  
À la conclusion de la fenêtre, ajouter \([A]_{10}\) à la liste des signatures.  
Conserver la position et la direction pour la fenêtre suivante.
\item Grouper les \(324\) indices par égalité de leurs listes de six
composants. Ces deux groupements donnent \(f_+\), \(f_\times\) et toutes
leurs préimages.
\item Pour chaque paire de signatures distinctes, appliquer \(G^{(6)}\).
Attribuer à l'émission la multiplicité produit des deux préimages.
Vérifier la récupération des signatures dans chaque bloc et la somme des
multiplicités \(104\,976\).
\item Appliquer \(\Psi\) à chaque émission et regrouper par égalité de mots.
Compter les \(243\) valeurs et les multiplicités de
\eqref{act21:tpk:fibras-ternarias}.
\item Appliquer \(r,s\) à chaque émission et vérifier son appartenance au catalogue.
Pour chaque mot ternaire, former ses six images diédriques, éliminer les
répétitions et enregistrer son minimum lexicographique. Compter les \(43\)
orbites et les deux tailles du tableau~\ref{act21:tpk:representantes}.
\end{enumerate}

L'exécution se termine : le premier tronçon évalue \(648\) trajectoires de \(54\) transitions ; le deuxième combine \(468\) couples de signatures ; le troisième regroupe un ensemble de \(243\) mots. Une vérification directe supplémentaire peut exécuter les \(104\,976\) couples et vérifier la factorisation \eqref{act21:tpk:factorizacion} ; sa portée coïncide avec le même domaine fini.

\subsection{Conservation du registre et fonction des germes}
\label{act21:tpk:conclusion}

Soit \(\mathsf{Sample}(q_t)\) l'échantillon antérieur à la transition et \(\mathcal M_t\) le registre accumulé. La mise à jour de la mémoire est
\begin{equation}
 \mathcal M_{t+1}=\mathcal M_t\mathbin{\Vert}\mathsf{Sample}(q_t),
 \label{act21:tpk:archivo}
\end{equation}
où \(\Vert\) désigne la concaténation ordonnée. La récurrence donne
\[
 \mathcal M_n=\mathcal M_0\mathbin{\Vert}
 (\mathsf{Sample}(q_0),\ldots,\mathsf{Sample}(q_{n-1})).
\]
Les accumulateurs se reconstruisent en additionnant les contributions actives de
chaque segment de neuf enregistrements. Leurs quotients, les tours des
curseurs et les orientations restent associés aux mêmes transitions.
L'exécution détermine ce registre au moyen de la récurrence ; le fichier
préserve la provenance concrète des lectures.

% BEGIN R27_MEM01
\subsection{Corrélation cyclique et conservation de l'ordre de la route}
\label{sec:rec27-memoria}

L'archive \eqref{act21:tpk:archivo} conserve la concaténation ordonnée des échantillons. L'exemple suivant précise une information d'ordre que deux lecteurs distincts extraient d'un mot de douze positions, en conservant le domaine de chaque lecture.

Soit $\tau=(\tau_1,\ldots,\tau_{12})\in\{-1,0,1\}^{12}$ un mot
ordonné, dont les indices sont interprétés modulo $12$. Nous définissons
\[
\begin{aligned}
 G_a&=\{a,a+4,a+8\}\quad(1\leq a\leq4),\\
 \nu_{120}(\tau)&=\sum_{a=1}^4\operatorname{sgn}
 \left(\sum_{m\in G_a}\tau_m\right),\\
 \nu_{270}(\tau)&=\sum_{m=1}^{12}\tau_m\tau_{m+3}.
\end{aligned}
\]
L'histogramme marginal est $p_j(\tau)=\#\{m:\tau_m=j\}/12$ pour
$j\in\{-1,0,1\}$ et son entropie est
$H(\tau)=-\sum_j p_j(\tau)\log p_j(\tau)$, avec $0\log0=0$.

\begin{proposition}[Information d'ordre distincte de la distribution marginale]
Il existe des mots ayant un histogramme identique et une corrélation différente
$\nu_{270}$, même en maintenant la même valeur de $\nu_{120}$.
\end{proposition}
\begin{proof}
Considérons
\[
 \tau^{(a)}=(1,1,0,0,0,0,0,0,0,0,0,0),\qquad
 \tau^{(b)}=(1,0,0,1,0,0,0,0,0,0,0,0).
\]
Toutes deux ont $p_1=1/6$, $p_0=5/6$ et $p_{-1}=0$, de sorte que leurs
entropies marginales coïncident. Dans chaque cas, deux des ensembles $G_a$
contiennent une unité et les deux autres uniquement des zéros : $\nu_{120}=2$.
Dans $\tau^{(a)}$, aucune paire d'unités n'a une séparation cyclique de trois,
tandis que dans $\tau^{(b)}$, seul le terme $m=1$ contribue.
Par conséquent, $\nu_{270}(\tau^{(a)})=0$ et $\nu_{270}(\tau^{(b)})=1$.
\end{proof}

Cette séparation identifie l'information conservée par le registre ordonné : les incidences entre positions, en plus des fréquences des symboles. Un état réduit qui ne retient que l'histogramme attribue la même sortie aux deux mots et perd cette corrélation. L'entropie marginale, l'incertitude entre bases et l'entropie d'intrication ont des objets de définition distincts ; leur articulation exige de conserver respectivement la distribution, la transformation de base et l'état avec sa partition.

% END R27_MEM01

% BEGIN R28_MEMORY_PARITY_INPUT
% BEGIN R28_MEMORY_PARITY
% Continuación aditiva de MEM01: los lectores ya están definidos.
\begin{proposition}[Parité des lecteurs sous inversion orientée]
Avec des indices cycliques dans \(\mathbb Z/12\mathbb Z\), soient
\((C_M\tau)_m=-\tau_{-m}\) et \(\operatorname{sgn}(0)=0\).
Les deux lecteurs satisfont
\[
 \nu_{120}(C_M\tau)=-\nu_{120}(\tau),\qquad
 \nu_{270}(C_M\tau)=\nu_{270}(\tau).
\]
\end{proposition}
\begin{proof}
La réflexion \(m\mapsto-m\) permute les quatre classes de positions
modulo quatre. L'inversion change le signe de chaque somme de classe et,
par l'imparité de \(\operatorname{sgn}\), celui de \(\nu_{120}\).
Dans le second lecteur, les signes des deux facteurs se compensent :
\[
 \sum_m(C_M\tau)_m(C_M\tau)_{m+3}
 =\sum_m\tau_{-m}\tau_{-m-3}
 =\sum_k\tau_{k+3}\tau_k
 =\nu_{270}(\tau).
\]
\end{proof}
L'orientation distingue ainsi une lecture impaire et une corrélation paire.
Toutes deux dépendent de la disposition des échantillons dans le registre ; leurs
lois de transformation précisent quelle information persiste lorsqu'on inverse
la route.
% END R28_MEMORY_PARITY

% END R28_MEMORY_PARITY_INPUT

À la fin de cette étape sont construits le domaine complet des conditions
initiales, l'émission décimale, ses deux signatures récupérables, l'image ternaire,
ses fibres et sa classification orbitale. Les recensements \(104\,976\), \(468\),
\(243\), \(729\) et \(43\) ont désormais un domaine et un mécanisme d'obtention.
Les sélections régionales recevront ces structures avec leurs invariants,
et les élévations agiront sur des mots orientés dont la provenance est
conservée. Cette continuité fait du catalogue fini un antécédent
constructif du développement coinductif.
